% Supplement: https://cp-algorithms.com/graph/hld.html
\subsection{树剖求 LCA}
\label{tree:hld-lca}
\AlgorithmInfo
    {利用重链剖分查询最近公共祖先，可与路径区间拆分共用预处理。}
    {预处理 $O(n)$，每次查询 $O(\log n)$，空间 $O(n)$。}
    {非空连通树，指定根；复用第~\ref{tree:hld}~节的 \InlineCode{HLD} 封装。}
\subsubsection*{原理}
重链剖分为每个点记录父亲 $fa$、深度 $dep$ 和所在链的链顶 $top$。查询 $u,v$ 时，若链顶不同，总让\textbf{链顶更深}的一侧跳到该链顶的父亲，即令 $u=fa_{top_u}$；若另一侧的链顶更深，先交换 $u,v$。

这一步不会越过 LCA：若较深的链顶是另一点的祖先，那么另一点所在链的链顶只能与它相同，或比它更深，均与当前选择矛盾。因此可以跳过当前点到链顶的整段。每次跳跃越过一条轻边，而每点到根至多有 $O(\log n)$ 条轻边，故查询复杂度为 $O(\log n)$。

两点进入同一条链后，深度较小的点就是 LCA。注意比较的是\textbf{链顶深度}，不能只比较当前点深度；同链时直接返回较浅点，也不必继续上跳。

倍增法适合同时跳祖先，树剖适合同时拆路径区间，两者 LCA 查询均为 $O(\log n)$。大量静态查询可用第~\ref{tree:rmq-lca}~节欧拉序 RMQ，单次 $O(1)$。

\subsubsection*{使用教程}
复制第~\ref{tree:hld}~节的 \InlineCode{HLD}，依次调用 \InlineCode{init(n)}、\InlineCode{add(u,v)} 和 \InlineCode{build(root)}，其中每条无向边加一次、根默认 1；随后 \InlineCode{lca(u,v)} 返回最近公共祖先。预处理不依赖线段树，查询支持两点相同或一方为另一方祖先。
\begin{lstlisting}[style=mtycpp,numbers=none]
HLD g;
g.init(5);
g.add(1,2),g.add(1,3),g.add(2,4),g.add(2,5);
g.build(1);
int p=g.lca(4,3);
int q=g.lca(2,4);
\end{lstlisting}
\subsubsection*{模板代码}
\InlineCode{HLD::lca(u,v)} 与两次 DFS 的完整实现统一放在第~\ref{tree:hld}~节。
